%% ut-thesis.tex -- document template for graduate theses at UofT
%%
%% Copyright (c) 1998-2013 Francois Pitt <fpitt@cs.utoronto.ca>
%% last updated at 16:20 (EDT) on Wed 25 Sep 2013
%%
%% This work may be distributed and/or modified under the conditions of
%% the LaTeX Project Public License, either version 1.3c of this license
%% or (at your option) any later version.
%% The latest version of this license is in
%%     http://www.latex-project.org/lppl.txt
%% and version 1.3c or later is part of all distributions of LaTeX
%% version 2005/12/01 or later.
%%
%% This work has the LPPL maintenance status "maintained".
%%
%% The Current Maintainer of this work is
%% Francois Pitt <fpitt@cs.utoronto.ca>.
%%
%% This work consists of the files listed in the accompanying README.

%% SUMMARY OF FEATURES:
%%
%% All environments, commands, and options provided by the `ut-thesis'
%% class will be described below, at the point where they should appear
%% in the document.  See the file `ut-thesis.cls' for more details.
%%
%% To explicitly set the pagestyle of any blank page inserted with
%% \cleardoublepage, use one of \clearemptydoublepage,
%% \clearplaindoublepage, \clearthesisdoublepage, or
%% \clearstandarddoublepage (to use the style currently in effect).
%%
%% For single-spaced quotes or quotations, use the `longquote' and
%% `longquotation' environments.


%%%%%%%%%%%%         PREAMBLE         %%%%%%%%%%%%

%%  - Default settings format a final copy (single-sided, normal
%%    margins, one-and-a-half-spaced with single-spaced notes).
%%  - For a rough copy (double-sided, normal margins, double-spaced,
%%    with the word "DRAFT" printed at each corner of every page), use
%%    the `draft' option.
%%  - The default global line spacing can be changed with one of the
%%    options `singlespaced', `onehalfspaced', or `doublespaced'.
%%  - Footnotes and marginal notes are all single-spaced by default, but
%%    can be made to have the same spacing as the rest of the document
%%    by using the option `standardspacednotes'.
%%  - The size of the margins can be changed with one of the options:
%%     . `narrowmargins' (1 1/4" left, 3/4" others),
%%     . `normalmargins' (1 1/4" left, 1" others),
%%     . `widemargins' (1 1/4" all),
%%     . `extrawidemargins' (1 1/2" all).
%%  - The pagestyle of "cleared" pages (empty pages inserted in
%%    two-sided documents to put the next page on the right-hand side)
%%    can be set with one of the options `cleardoublepagestyleempty',
%%    `cleardoublepagestyleplain', or `cleardoublepagestylestandard'.
%%  - Any other standard option for the `report' document class can be
%%    used to override the default or draft settings (such as `10pt',
%%    `11pt', `12pt'), and standard LaTeX packages can be used to
%%    further customize the layout and/or formatting of the document.

%% *** Add any desired options. ***


\documentclass[11pt]{ut-thesis}

\usepackage{appendix}
\usepackage{graphicx}
\usepackage{amsmath}
\usepackage{amssymb}
\usepackage{amsthm}
\usepackage{mathrsfs}
\usepackage[byname]{smartref}
\usepackage{hyperref} %comment out for hardcopy
%\usepackage{txfonts}
\usepackage{tocloft}
\usepackage[numbers]{natbib}
\usepackage[symbols,nogroupskip,nonumberlist]{glossaries-extra}

\makeglossaries

\glsxtrnewsymbol[description={Individual index}]{i}{\ensuremath{i}}
\glsxtrnewsymbol[description={Family (Cluster) index}]{j}{\ensuremath{j}}
\glsxtrnewsymbol[description={Proband index}]{p}{\ensuremath{p}}
\glsxtrnewsymbol[description={Number of events in family $j$}]{dj}{$d_j$}
\glsxtrnewsymbol[description={Some time}]{t}{\ensuremath{t}}
\glsxtrnewsymbol[description={Some Time for the proband}]{a}{\ensuremath{a}}
\glsxtrnewsymbol[description={Event Time}]{T}{\ensuremath{T}}
\glsxtrnewsymbol[description={Event indicator for individual $i$ in family $j$}]{delta}{$\delta_{ij}$}
\glsxtrnewsymbol[description={The observed survival data ($t,\delta$)}]{w}{\ensuremath{w}}
\glsxtrnewsymbol[description={Number of individuals}]{n}{\ensuremath{n}}
\glsxtrnewsymbol[description={Number of Families (Clusters)}]{J}{\ensuremath{J}}
\glsxtrnewsymbol[description={Index of the sampled completed dataset in the MCEM}]{m}{\ensuremath{m}}
\glsxtrnewsymbol[description={Number of the sampled completed dataset in the MCEM}]{M}{\ensuremath{M}}
\glsxtrnewsymbol[description={Frailty term}]{z}{\ensuremath{z}}
\glsxtrnewsymbol[description={$q$-th element of Gauss Hermite Quadrature}]{q}{\ensuremath{q}}
\glsxtrnewsymbol[description={$q$-th weight of Gauss Hermite Quadrature}]{omega}{\ensuremath{\omega}}
\glsxtrnewsymbol[description={$q$-th node of Gauss Hermite Quadrature}]{yq}{\ensuremath{y_q}}
\glsxtrnewsymbol[description={Total number of quadratures}]{Nq}{\ensuremath{N_q}}
\glsxtrnewsymbol[description={Hazard fucntion}]{h}{\ensuremath{h(\cdot)}}
\glsxtrnewsymbol[description={Baseline hazard function}]{h_0}{\ensuremath{h_0(\cdot)}}
\glsxtrnewsymbol[description={Cumulative hazard fucntion}]{H}{\ensuremath{H(\cdot)}}
\glsxtrnewsymbol[description={Survival fucntion}]{S}{\ensuremath{S(\cdot)}}
\glsxtrnewsymbol[description={Ascertainment of family $j$ into the study}]{A}{$A_j(\cdot)$}
\glsxtrnewsymbol[description={Likelihood function}]{L}{$L(\cdot)$}
\glsxtrnewsymbol[description={Log-likelihood function}]{ell}{$\ell(\cdot)$}
\glsxtrnewsymbol[description={Laplace transform}]{lap}{$\mathscr{L}(\cdot)$}
\glsxtrnewsymbol[description={Covariates}]{X}{\ensuremath{\mathbf{x}}}
\glsxtrnewsymbol[description={Model coefficients vector}]{beta}{$\boldsymbol{\beta}$}
\glsxtrnewsymbol[description={Parameter vector}]{theta}{$\boldsymbol{\theta}$}
\glsxtrnewsymbol[description={The combination of $(\boldsymbol{\beta}, \lambda, \alpha)$}]{Lambda}{$\Lambda$}
\glsxtrnewsymbol[description={Weibull shape parameter}]{lambda}{$\lambda$}
\glsxtrnewsymbol[description={Weibull scale parameter}]{alpha}{$\alpha$}
\glsxtrnewsymbol[description={General form of the parameter in an undefined frailty distribution}]{upsilon}{$\upsilon$}
\glsxtrnewsymbol[description={Gamma shape and rate parameters}]{k}{$k$}
\glsxtrnewsymbol[description={Log-Normal variance parameter}]{sigma}{$\sigma^2$}
\glsxtrnewsymbol[description={Missing data distribution parameters}]{psi}{$\psi$}



%% *** Add \usepackage declarations here. ***
%% The standard packages `geometry' and `setspace' are already loaded by
%% `ut-thesis' -- see their documentation for details of the features
%% they provide.  In particular, you may use the \geometry command here
%% to adjust the margins if none of the ut-thesis options are suitable
%% (see the `geometry' package for details).  You may also use the
%% \setstretch command to set the line spacing to a value other than
%% single, one-and-a-half, or double spaced (see the `setspace' package
%% for details).


%%%%%%%%%%%%%%%%%%%%%%%%%%%%%%%%%%%%%%%%%%%%%%%%%%%%%%%%%%%%%%%%%%%%%%%%
%%                                                                    %%
%%                   ***   I M P O R T A N T   ***                    %%
%%                                                                    %%
%%  Fill in the following fields with the required information:       %%
%%   - \degree{...}       name of the degree obtained                 %%
%%   - \department{...}   name of the graduate department             %%
%%   - \gradyear{...}     year of graduation                          %%
%%   - \author{...}       name of the author                          %%
%%   - \title{...}        title of the thesis                         %%
%%%%%%%%%%%%%%%%%%%%%%%%%%%%%%%%%%%%%%%%%%%%%%%%%%%%%%%%%%%%%%%%%%%%%%%%

%% *** Change this example to appropriate values. ***
\degree{Master of Science}
\department{Epidemiology and Biostatistics\\
Schulich School of Medicine \& Dentistry}
\gradyear{2024}
\author{Jiaqi Bi}
\title{Frailty Model Incorporating Ascertainment Correction with Missing Data in Family-Based Study}

%% *** NOTE ***
%% Put here all other formatting commands that belong in the preamble.
%% In particular, you should put all of your \newcommand's,
%% \newenvironment's, \newtheorem's, etc. (in other words, all the
%% global definitions that you will need throughout your thesis) in a
%% separate file and use "\input{filename}" to input it here.


%% *** Adjust the following settings as desired. ***

%% List only down to subsections in the table of contents;
%% 0=chapter, 1=section, 2=subsection, 3=subsubsection, etc.
\setcounter{tocdepth}{2}

%% Make each page fill up the entire page.
\flushbottom


%%%%%%%%%%%%      MAIN  DOCUMENT      %%%%%%%%%%%%

\begin{document}

%% This sets the page style and numbering for preliminary sections.
\begin{preliminary}

%% This generates the title page from the information given above.
\maketitle

%% There should be NOTHING between the title page and abstract.
%% However, if your document is two-sided and you want the abstract
%% _not_ to appear on the back of the title page, then uncomment the
%% following line.
\cleardoublepage

%% This generates the abstract page, with the line spacing adjusted
%% according to SGS guidelines.
\begin{abstract}
    \addcontentsline{toc}{chapter}{Abstract}
    This is an abstract section...
%% *** Put your Abstract here. ***
%% (At most 150 words for M.Sc. or 350 words for Ph.D.)
\end{abstract}

%% Anything placed between the abstract and table of contents will
%% appear on a separate page since the abstract ends with \newpage and
%% the table of contents starts with \clearpage.  Use \cleardoublepage
%% for anything that you want to appear on a right-hand page.

%% This generates a "dedication" section, if needed -- just a paragraph
%% formatted flush right (uncomment to have it appear in the document).
\begin{dedication}
    \addcontentsline{toc}{chapter}{Dedication}
    \textit{To my mother, Kai Hua, for all the upbringrings. }

    \textit{To my father, Guangmin Bi, for teaching me not to give up. }
    
    \textit{To my supervisors, for their unstopping guidance and invaluable lessons. }
%% *** Put your Dedication here. ***
\end{dedication}

%% The `dedication' and `acknowledgements' sections do not create new
%% pages so if you want the two sections to appear on separate pages,
%% uncomment the following line.
\newpage  % separate pages for dedication and acknowledgements

%% Alternatively, if you leave both on the same page, it is probably a
%% good idea to add a bit of extra vertical space in between the two --
%% for example, as follows (adjust as desired).
%\vspace{.5in}  % vertical space between dedication and acknowledgements

%% This generates an "acknowledgements" section, if needed
%% (uncomment to have it appear in the document).
\begin{acknowledgements}
    \addcontentsline{toc}{chapter}{Acknowledgements}
%% *** Put your Acknowledgements here. ***
\end{acknowledgements}

\newpage

%% This generates the Table of Contents (on a separate page).
\tableofcontents
\printunsrtglossary[type=symbols,style=long,title={List of Symbols}]

\newpage
%% This generates the List of Tables (on a separate page), if needed
%% (uncomment to have it appear in the document).
\listoftables

\newpage
%% This generates the List of Figures (on a separate page), if needed
%% (uncomment to have it appear in the document).
\listoffigures

%% You can add commands here to generate any other material that belongs
%% in the head matter (for example, List of Plates, Index of Symbols, or
%% List of Appendices).

%% End of the preliminary sections: reset page style and numbering.
\end{preliminary}


%%%%%%%%%%%%%%%%%%%%%%%%%%%%%%%%%%%%%%%%%%%%%%%%%%%%%%%%%%%%%%%%%%%%%%%%
%%  Put your Chapters here; the easiest way to do this is to keep     %%
%%  each chapter in a separate file and `\include' all the files.     %%
%%  Each chapter file should start with "\chapter{ChapterName}".      %%
%%  Note that using `\include' instead of `\input' will make each     %%
%%  chapter start on a new page, and allow you to format only parts   %%
%%  of your thesis at a time by using `\includeonly'.                 %%
%%%%%%%%%%%%%%%%%%%%%%%%%%%%%%%%%%%%%%%%%%%%%%%%%%%%%%%%%%%%%%%%%%%%%%%%
\chapter{Time series: Long memory}
Here is a picture of a long memory time series. 
\begin{figure}[ht]
\begin{center}
%\includegraphics[height = 9cm, width = 9cm]{pic1.jpeg}
\includegraphics[height = 9cm, width = 9cm]{logo.png}
\caption{A long memory time series\label{ts1}}
\end{center}
\end{figure}

Here's a table.
\begin{table}[ht]
\begin{center}
\begin{tabular}[ht]{|c|lr|c|} 
%c stands for centre, l for left, r for right; the | puts lines in between, and the hline puts a horizontal line in
\hline
$n$ & $\alpha$ &$n\alpha$ & $\beta$\\
\hline
1 & 0.2 & 0.2 & 5\\
\hline
2 & 0.3 & 0.6 & 4\\
\hline
3 & 0.7 & 2.1 & 3\\
\hline
\end{tabular}
\caption{A random table \label{tab1}}
\end{center}
\end{table}

\begin{eqnarray}
y &=& mx + b \label{eq1}\\
&=& ax+ c
\label{eq2}
\end{eqnarray}

This is an un-numbered equation, along with a numbered one. 
\begin{eqnarray}
u &=& px \nonumber\\
p &=& P(X=x) \label{eqn3}
\end{eqnarray}

Look at Table \ref{tab1} and Figure \ref{ts1} and equations \ref{eq1},  \ref{eq2}, and \ref{eqn3}.

Let's do some matrix algebra now.

\begin{equation}
det\left(\left|\begin{array}{ccc} 2 & 3 & 5\\
4 & 4 & 6\\
9 & 8 & 1
\end{array}\right|\right) = 42
\end{equation}

In the equation and eqnarray environments, you don't need to have the dollar sign to enter math mode.

\begin{eqnarray}
\alpha = \beta_1 \Gamma^{-1}
\end{eqnarray}

This is citing a reference ~\cite{mygood11111}.  This is citing another ~\cite{mrx05}.  Nobody said something ~\cite{Nobody06}.

%% *** Include chapter files here. ***


%% This adds a line for the Bibliography in the Table of Contents.
\addcontentsline{toc}{chapter}{Bibliography}
%% *** Set the bibliography style. ***
%% (change according to your preference/requirements)
\bibliographystyle{unsrtnat}
%% *** Set the bibliography file. ***
%% ("thesis.bib" by default; change as needed)
\bibliography{westernthesis.bib}

%% *** NOTE ***
%% If you don't use bibliography files, comment out the previous line
%% and use \begin{thebibliography}...\end{thebibliography}.  (In that
%% case, you should probably put the bibliography in a separate file and
%% `\include' or `\input' it here).

\end{document}
